\section{Boundary obstructions, links, and strict gap growth}
\label{sec:links}

An escaping source translation acts on a finite regular array. We will
record the limiting flags of its measured Fourier index and compare the
measurements seen from the two ends of the translation. The first task is
to construct these measurements and their input forms. The second is to
prove that a vertex translation produces exactly one divergent target
gap, and that successively completing such a vertex-origin path
adds exactly one target gap at each stage.

We retain the regular realizations and induced actions of
Section~\ref{sec:setup}. In particular,
\(H=eL^2(B)^k\), \(\mathcal H=L^2(Y;H)\), and
\(d=(\Tr\otimes\tau_B)(e)\). For an integrable corner field write
\[
 \mathsf T(C)=\int_Y(\Tr\otimes\tau_B)(C(y))\,dy,
 \qquad \tau_{\mathcal N}(C)=d^{-1}\mathsf T(C).
\]
Thus \(\tau_{\mathcal N}\) is a probability trace, whereas the Hilbert
norm of a column satisfies \(\|f\|_2^2=\mathsf T(ff^*)\).
All matrix sizes in this section are fixed and finite. An
operator-bounded column means an essentially bounded field of bounded
tracial columns. Such columns form a dense subspace of \(\mathcal H\).
The constant corner frame will be denoted by
\[
 b_r=e\varepsilon_r,\qquad 1\leq r\leq k,
 \qquad \sum_{r=1}^k b_rb_r^*=e.
\]
Here and below a constant column can still have entries in \(B\).

\subsection{Compact matrix data and the surviving flags}

Write \(\Pi\) for the indivisible relative simple roots, over all
components. A face of type \(I\subset\Pi\) has strictly positive
simple coordinates exactly in \(I\); opposition of types is denoted
by a star. The Cartan projection is written \(\mu(h)\) in this
geometric discussion, and is distinct from a scalar cocycle. Partial
flags are transverse when they are in the open relative position of
subfaces of opposite chambers. Incidence means containment in one
chamber.

\begin{lemma}[Matrix compactification and residual projectivities]
\label{lem:links-matrix-data}
There is a finite family of projectively normalized matrix coordinates,
including coordinates of the inverse, with the following properties.
\begin{enumerate}[label=(\roman*)]
\item The set \(D\subset\Pi\) of diverging simple gaps is detected
by rank one in the fundamental highest-line coordinates. A stratum
with gap set \(D\) carries an image \(D\)-flag, a reverse
\(D^*\)-flag, and a residual projectivity between their Levi
refinement spaces. Its closure has only gap sets containing \(D\).
\item On a bounded complementary-gap window, the residual blocks are
invertible in compact frames. They send a full flag transverse to the
reverse face, by transverse parabolic projection followed by Levi
transport, to a full flag incident with the image face.
\item Fixed left and right group multiplications preserve gap sets.
Right multiplication leaves the image flag unchanged; left
multiplication transports it by the usual target action. Reverse data
obey the opposite rule, and residual projectivities change by the
corresponding compositions.
\item In an archimedean highest-line module, each surviving top face
space is spanned by the highest lines of flags incident with the
observed face. A rank-one Levi pencil has the effective rank-one
projectivity group and its omitted simple gap, up to the bounded
changes made by compact endpoint frames.
\end{enumerate}
The matrix coordinates are real smooth coordinates at archimedean
places and projective local-field coordinates at finite places.
\end{lemma}

\begin{proof}
Use the highest-line modules, their normalized matrices and their inverses
from Section~\ref{sec:setup}. At archimedean places normalize by the
Hilbert--Schmidt norm; at finite places use projective ratio charts
with a maximal coordinate. Lemma~\ref{lem:setup-face-blocks} identifies
the surviving Levi blocks when the gap set is $D$, proves their
invertibility on a bounded complementary-gap window, and describes the
intrinsic transverse projection followed by residual transport. In each
fundamental highest-line coordinate, rank one detects divergence of the
corresponding simple gap. These finitely many coordinates therefore
record the gap set, its image and reverse flags, and the residual maps.
If an additional gap diverges, the previous image vertices persist by
that lemma. This proves the stratum and closure assertions in
\textup{(i)}, and the residual-map assertion in \textup{(ii)}.

We verify the fixed-multiplication rule that will be used on Fourier
indices. For a fixed matrix $a$, both $a$ and $a^{-1}$ have finite norm.
Thus projectively normalized multiplication by $a$ has a denominator
bounded away from zero. Applying this to every chosen module and its
inverse shows that fixed multiplication changes each simple Cartan gap
by a bounded amount. Left multiplication carries the image top space
by $a$, whereas right multiplication leaves it unchanged. Applying the
same calculation to inverse matrices gives the reverse-flag rule.
The residual map transforms by the corresponding compositions, since
its description by transverse projection and surviving Levi transport
is intrinsic. This proves \textup{(iii)}.

The last assertion is the spanning and length statement of
Lemma~\ref{lem:setup-face-blocks}: the Levi orbit of the highest line
spans its surviving top space, and on a panel the only noncompact
isotropic factor supplies the omitted simple gap. Compact changes of
endpoint frame change that length by a bounded amount. At a finite
place the modules remain vector spaces over the local field; no
Hilbert-space amplification by them is involved.
\end{proof}

Compact windows in a stratum are exhausted by bounds on its
complementary gaps and by compact transverse projective charts.
Their tests may be chosen with slack: a slightly larger window
retains the same large-gap separation and the same lower bounds for
the invertible residual blocks. At real places these tests are
smooth functions of matrix coordinates. At finite places they are
locally constant on compact boxes at a fixed level. Increasing a
highest weight by tensor powers in later polynomial tests does not
change the scalar projective multiplier of the original Hilbert
representation.

\subsection{Fourier transport before compression}

The Fourier tests in this section act on the full regular arrays.
They are not followed by a corner projection before the next
operation. This distinction allows us to use a diagonal spectral
calculus without claiming that an angular Fourier cutoff preserves
the operator norm of a column.

Choose a path \(a\) with \(a^{-1}\) split toward a source face
\(F\) of type \(I\), with its coordinates in \(I\) tending to
infinity. Conjugate split frames are allowed. A \emph{lexicographic
face path} starts with a vertex translation and adds one new simple
coordinate at each successively slower scale. Its fastest parameter
is taken to its limit first, while all slower parameters are fixed;
the next limit is then taken, and so on. A general face path may
instead have several coordinates tending to infinity simultaneously.
The general gap upper bound below applies to these paths as well;
its exact successive count is asserted for the specified
vertex-origin chains. Finite-place parameters range on the divisible
cocharacter meshes fixed in Section~\ref{sec:setup}.

The measured label is the compact matrix data of $h^{-1}$, where $h$
is the Fourier index in the chosen copy $\Gamma_j$. The other group
indices remain coefficient indices. On the full regular arrays, a scalar
test of this label is a diagonal Hilbert-space multiplier. These tests
are not followed by corner compression before another array operation.

There are two input maps. The one-column map is
\[
 (V_a f)(y)=X_a(y)f(y).
\]
It is an isometry with the integrated tracial Hilbert norm. Its Hilbert
ultralimit $V$ is therefore an isometry. On the minimal spectral span
generated by $V\mathcal H$, the continuous diagonal tests give a
projection-valued measure $E^{(1)}$. Its compression
\[
 \Phi(A)=V^*E^{(1)}(A)V
\]
is a positive effect on the original input. When $A$ records only a gap
stratum and its image flag, this is the image measurement. It need not
be a projection.

The two-column map, initially on operator-bounded columns, is
\begin{equation}
 J(\overline b,f)
  =\sqrt d\,\bigl(b(ay)^*X_a(y)f(y)\bigr)^{\lim}.
 \label{eq:input-1}
\end{equation}
Its output has its own diagonal spectral measure, denoted by $E$.
The link identities below determine the norm of this map and the
Hilbert inputs to which it extends. Thus the symbol $J$ at this point
specifies a map on bounded column pairs, not a tensor-product isometry.

In both constructions, iterated ultraproducts realize the specified
successive limits. Restrict to the minimal closed spectral span generated
by the input ranges and a countable determining algebra of continuous
tests; these spaces are separable. A Borel stratum always denotes a
projection of the limiting spectral measure, rather than a discontinuous
cutoff applied pointwise before taking the limit.

\begin{lemma}[Fourier--Gram transport]
\label{lem:links-gram-transport}
On the limiting arrays, a fixed left or right Fourier shift
transports continuous compact matrix tests by its limiting geometric
map, and hence transports the associated Borel spectral projections.
For a stratum and its image flag, a left Fourier matrix multiplier
can be removed from a quadratic form through its uncut Gram factor.
This remains true for bounded base-dependent multipliers and for
compact families of fixed source translations reindexed at the end
base. All projective defects cancel in these sesquilinear identities.
\end{lemma}

\begin{proof}
For a fixed shift, normalized matrix multiplication and the inverse
formula have denominators bounded away from zero by
Lemma~\ref{lem:links-matrix-data}. Their action on the compact label
space is continuous. A diagonal continuous test therefore
intertwines the shift in the ultralimit. Equality of the two
representations on continuous functions implies equality of their
spectral measures, so the assertion extends to Borel projections of
the limiting variables. Base multiplication and the reindexing of
the base commute with diagonal Fourier tests.

If a Fourier index \(h\) is changed to \(rh\) by a left
multiplier, the recorded inverse is changed to \(h^{-1}r^{-1}\).
Thus its gap set and image flag are unchanged in the limit. Let
\(D\) be a bounded test of those data and \(T\) a bounded left
Fourier matrix multiplier. Choose uniformly bounded finite Fourier
approximants \(T_l\to T\) in trace \(L^2\). In a form tested on
uncut operator-bounded columns, replacing either occurrence of
\(T\) by \(T_l\) before moving \(D\) has error bounded by
\[
 C\|D\|\,\|T-T_l\|_2,
\]
where \(C\) depends only on the fixed column bounds and matrix
sizes. For each finite Fourier sum, move \(D\) through its
individual shifts by the already proved limiting identity. The
remaining product is the uncut Gram factor \(T_l^*T_l\). Replacing
this by \(T^*T\) costs another quantity tending to zero, since
\[
 \|T_l^*T_l-T^*T\|_2
 \leq (\|T_l\|_\infty+\|T\|_\infty)\|T_l-T\|_2.
\]
All these errors are estimated before applying the cutoff to the
columns. In particular the argument applies to a partial unitary
with \(T^*T=e\) and does not assign an operator bound to \(Df\).

For a field reindexed at \(ay\), invariant probability on \(Y\)
gives
\(\int\|T(ay)-T_l(ay)\|_2^2dy=\|T-T_l\|_2^2\).
The same estimates therefore apply after path transport. For a
compact family of fixed translations, the cocycle formula and
continuity of translations in measure make their matrix fields a
compact family in trace \(L^2\). To verify this assertion directly,
approximate the Borel section by continuous data on sets of large
measure; along a convergent sequence of translations the discrete
cocycle labels then agree eventually off sets of small measure.
Unitarity bounds the remaining error. A finite net gives uniform
finite-Fourier approximation of the compact family. No common
finite-Fourier approximation of the escaping paths themselves is
used.

Finally, the group-valued section cocycle satisfies its exact
composition identity. Applying the projective source representation
gives only a scalar function of the base in each composition
formula. Pulling a fixed source translation from one end of an
array to the other thus multiplies the \emph{whole} array by one
scalar of modulus one. It commutes with diagonal Fourier tests and
cancels between the two entries of every sesquilinear form. The
same observation applies to a quotient-path calculation: after
base reindexing, \(X_tX_a^*\) is the transporter for \(ta^{-1}\)
times a common scalar, so its localized Fourier norm is unchanged.
Products in these identities are the actual twisted matrix
products; their coefficient phases cannot be replaced by absolute
values of untwisted convolutions.
\end{proof}

When there is one measured coordinate, the other group indices,
finite matrix factors, and base algebra may be left in the
coefficient algebra. The multiplier splits across these group
factors, so the proof just given applies without an additional
commutation assumption on coefficient entries.

\subsection{Link forms, scalar masses, and the finite frame}

We first compute the input norm and the covariances of the measurements.
For a stratum of positive scalar mass $m_D$, divide its image measurement
by $m_D$ to obtain a unital measurement. Scalarity is part of the next
proposition, rather than an assumption on the construction.

\begin{proposition}[Link identities]
\label{prop:links-link}
The one-column measurements and the forms
\eqref{eq:input-1} have the following properties.
\begin{enumerate}[label=(\roman*)]
\item The image measurement in one measured copy lies in
\(\mathcal A_j\), is \(\beta_j\)-equivariant, and is invariant
under the source face parabolic. Its gap-stratum masses belong to
\(\mathcal N\) and are scalar. Conjugating the source split frame
transports the measurement by \(\sigma\).
\item If the path escapes in every simple component, the link is
an isometry
\(\overline{\mathcal H}\otimes\mathcal H\to\mathcal K\).
Its output forms are decomposable at the two bases, covariant under
the two independent target shifts, invariant under the independently
contracted source radicals, and invariant under simultaneous source
centralizers. An image-and-stratum effect depends only on the start
input and is \(1\otimes\Phi(A)\); the reverse rule is symmetric.
\item If the path is active precisely in a subproduct \(G_0\),
put \(N_0=\mathcal N^{G_0}\), and let
\(E_{N_0}:\mathcal N\to N_0\) be the trace-preserving conditional
expectation. On bounded columns the link instead has the form
\begin{equation}
 \|J(\overline b,f)\|^2
  =d^2\tau_{\mathcal N}
       \bigl(E_{N_0}(ff^*)E_{N_0}(bb^*)\bigr).
 \label{eq:links-relative-form}
\end{equation}
Its polarized extension on finite sums is a positive semidefinite
form. The same Gram covariances hold on this form space; they do
not assert a tensor-product isometry. The precise Hilbert-input
extensions are given in Lemma~\ref{lem:links-relative-boundedness}.
\item For any Borel event \(A\) determined by a gap stratum and
its image flag, the finite-frame identity is
\begin{equation}
 \langle f,\Phi(A)f\rangle
  =\frac1d\sum_{r=1}^k
       \|E(A)J(\overline{b_r},f)\|^2.
 \label{eq:links-frame}
\end{equation}
It holds on mixed and nonmixed paths. More generally it holds for
events in additional limiting coordinates which are unchanged by
the ignored left Fourier shifts. When a proof starts with windows
that also constrain the reverse flag or residual projectivity,
those windows must first be exhausted to such an invariant event.
\item The image measurements are independent of the finite regular
embedding. Adding another original tensor slot gives their normal
copies, tensored with the identity on that slot. This assertion
concerns these compressed measurements and operators determined by
their stated one-ended covariance and uniqueness.
\end{enumerate}
All statements use the specified order of successive limits. A
contracted translate must tend to the identity in that order.
\end{proposition}

\begin{proof}
\emph{(i) One-ended covariance.}
Base multipliers and the companion actions commute with the relevant
diagonal tests and intertwine the input column map. The compressed
image effects therefore lie in \(\mathcal A_j\). A measured target
shift changes the inverse Fourier label by the corresponding left
geometric action. Lemma~\ref{lem:links-gram-transport} consequently
gives \(\beta_j\)-equivariance, with the action on functions given
by composition with the inverse target element.

Let \(n\) be in the radical contracted by the path. Sliding
\(U_n\) through \(U_a\) puts \(U_{ana^{-1}}\) at the far end,
up to a common base phase. This translation tends strongly to the
identity. In the section model its discrete cocycle is the identity
off a set of measure tending to zero, after its pure base shift has
been undone. On bounded inputs both descriptions give a vanishing
Hilbert error, even with a bounded diagonal test inserted. This
proves radical invariance of the full limiting data.

For a fixed centralizing source element \(s\), the corresponding
identity slides \(U_s\) to the far end with the same \(s\).
Its base reindexing has no effect on an integrated form, and its
left corner multiplier can be removed for image-and-stratum tests
by Lemma~\ref{lem:links-gram-transport}. Its uncut Gram factor is
the corner identity. Thus the image measurement is centralizer
invariant. The radical and this Levi centralizer generate the face
parabolic. The same calculation for a conjugated split path removes
the end conjugation and leaves precisely source transport of its
image marginal.

The mass of a gap stratum is target invariant, hence belongs to
\(\mathcal N\). It is invariant under the entire source face
parabolic. Every inactive simple factor is in the centralizer. In
an active factor the contracted radical contains an escaping
sequence. Apply Howe--Moore to the representation of that factor
on \(L^2(\mathcal N)\): a radical-invariant vector orthogonal to
the factor-invariant subspace would have a matrix coefficient both
constant and tending to zero. It must therefore be factor invariant.
The mass is invariant under all simple factors and so is scalar,
because \(\mathcal N^G=\mathbb C\). This argument applies to a
pure vertex path as well as to a mixed path.

\emph{(ii) The mixed norm and its polarization.}
For bounded \(b,f\), tracial cyclicity and the definition of the
induced action give
\[
 \|J(\overline b,f)\|^2
  =d\lim_a\mathsf T\bigl(ff^*\sigma_{a^{-1}}(bb^*)\bigr).
\]
On a mixed path, trace mixing makes the last expression
\[
 d^2\tau_{\mathcal N}(ff^*)\tau_{\mathcal N}(bb^*)
       =\|f\|_2^2\|b\|_2^2.
\]
Polarization, or the same calculation on finite sums, gives
\[
 \left\|\sum_iJ(\overline{b_i},f_i)\right\|^2
  =d^2\sum_{i,l}
      \tau_{\mathcal N}(f_lf_i^*)
      \tau_{\mathcal N}(b_ib_l^*).
\]
This is exactly the norm on
\(\overline{\mathcal H}\otimes\mathcal H\). The factor
\(\sqrt d\) in \eqref{eq:input-1} is essential for this
normalization.

Multiplication of an input by a start base function is represented
on arrays by that function at \(y\); an end base function is
represented by its value at \(ay\). The diagonal tests commute
with both. Their limiting forms are therefore decomposable at the
two bases. Moving an independent target shift in either input
uses the left or right Fourier transport formula and proves the
two covariances. A source radical in either input is slid to its
contracted translate in the other input, as above. The fixed
translated input converges strongly, its column bound is preserved,
and the common phase cancels. Simultaneous centralizing translations
use the same exact precomposition and finite-Fourier Gram argument.

For an image-and-stratum event, the ignored end Fourier shifts put
each end slice of the compressed effect in
\(\overline{\mathcal N}\). Independent reverse-radical
invariance makes this slice scalar by the factorwise Howe--Moore
argument, since every factor is active. Thus the compressed effect
has the form \(1\otimes C\). The frame computation below
identifies \(C=\Phi(A)\). Reversing the two inputs proves the
reverse statement.

\emph{(iii) The nonmixed form.}
For an escape in \(G_0\), the weak limit of the trace-preserving
action on \(L^2(\mathcal N)\) is the projection onto
\(L^2(N_0)\). Repeating the norm computation gives
\eqref{eq:links-relative-form}. For a finite sum the formula is
\begin{equation}
 \left\|\sum_iJ(\overline{b_i},f_i)\right\|^2
 =d^2\sum_{i,l}\tau_{\mathcal N}
       \bigl(E_{N_0}(f_lf_i^*)E_{N_0}(b_ib_l^*)\bigr).
 \label{eq:links-polarized-relative}
\end{equation}
It is positive because it is the limit of squared norms. Quotient
by its null space and complete when a Hilbert form space is needed.
Every covariance proved before the mixed norm calculation remains
an equality of these Gram forms. In particular it can be used on
the relative-boundedness domains below. The mass calculation in
the first part of the proof is separate from this formula and
does not replace \(N_0\) by the scalars.

\emph{(iv) Frame recovery.}
For a continuous test of the image and stratum, approximate each
left multiplier \(b_r(ay)^*\) by bounded finite Fourier matrices.
The inverse-index convention makes the test invariant under their
left Fourier shifts. Lemma~\ref{lem:links-gram-transport} therefore
moves the test past these multipliers. Summing their uncut Gram
factors gives
\(\sum_r b_r(ay)b_r(ay)^*=e\). The two-column normalization
contributes exactly \(d\), proving
\eqref{eq:links-frame}. Equality of the limiting spectral measures
extends it to invariant Borel events. The proof used neither
factorization of the end variable nor mixed trace mixing.

For a window involving residual or reverse data the left shifts
may change the window, so that argument does not apply to the
individual window. Suppose instead that all end vectors have been
annihilated on each member of an exhaustion of the desired
image-and-stratum event. Monotone spectral convergence first gives
annihilation on that event, and only then may the finite frame be
inserted. This proves the stated order of operations. The same
proof covers extra limiting coordinates whenever they too are
unchanged by the ignored shifts.

\emph{(v) Embedding independence and normal copies.}
An intertwiner between two regular embeddings is a bounded Fourier
matrix on the relevant finite modules. For the image-and-stratum
tests, its uncut Gram factor is the appropriate input corner
identity. Lemma~\ref{lem:links-gram-transport} consequently
identifies the two compressed measurements.

For the normal-copy assertion, begin with a regular-multiple
embedding for the first two slots and enlarge it by the entire
extra slot. At the regular \(\Lambda\)-index \(h\), conjugate
the extra slot by the projective inverse of its simultaneous
\(\rho(h)\) action. Projective absorption transfers its
multiplier to the regular \(\Lambda\) multiplier. This conjugation
commutes with the extra \(\Gamma\)-action, which can now be
embedded in its own finite regular multiple. Both operations
commute with scalar Fourier cutoffs on the original measured
\(\Gamma_j\)-coordinate. In the original input coordinates,
the one-ended path acts only in the zeroth slot. Its compressed
one-ended tests are therefore the original tests tensored with
identity on the added slot. The embedding-independence just proved
identifies this realization with any subsequent one. This proves
the assertion for the measurements, and hence for the normal
algebras and uniquely determined operators constructed from them.
It makes no assertion that all uncompressed spectral arrays are
embedding independent.
\end{proof}

The two norm formulas give different input spaces. On an all-factor
path, $J$ extends to the ordinary Hilbert tensor product. On a partial
path its norm still contains $E_{N_0}(ff^*)$, which can be unbounded for
a Hilbert column. The useful replacement is to bound this conditional
expectation for one input. Such a bound lets the other input vary through
its whole Hilbert space and survives the image effects used later.

\begin{lemma}[Relative boundedness and Hilbert inputs]
\label{lem:links-relative-boundedness}
For a path supported in \(G_0\), call a column \(f\) relatively
bounded if the positive affiliated operator
\(E_{N_0}(ff^*)\) is bounded. Then:
\begin{enumerate}[label=(\roman*)]
\item for fixed relatively bounded \(f\), the link extends
continuously to every Hilbert end column \(b\), and conversely
with the inputs reversed;
\item the image measurements of paths in factors contained in
\(G_0\) commute with \(N_0\); their generated algebra, matrix
effects, and support projections preserve relative boundedness;
\item a closed \(N_0\)-invariant subspace has a dense relatively
bounded part, and a relatively bounded column has operator-bounded
approximants with the same bound on its conditional expectation;
\item the link covariances hold on relatively bounded pairs and
when one input is operator-bounded and the other is a Hilbert
column. A one-column gap inclusion valid for all inputs also holds
for bounded pairs in \eqref{eq:input-1}.
\end{enumerate}
\end{lemma}

\begin{proof}
Conditional expectation extends to positive \(L^1\) elements.
Put \(A=E_{N_0}(ff^*)\). If \(A\) is bounded,
\eqref{eq:links-relative-form} gives
\begin{equation}
 \|J(\overline b,f)\|^2
       \leq d\|A\|_\infty\|b\|_2^2.
 \label{eq:links-relative-bound}
\end{equation}
This proves the first continuous extension from bounded \(b\).
The construction always takes the link limit before this extension;
it does not define a finite-time product of two arbitrary Hilbert
columns.

Let \(c\in N_0\). Along a path in \(G_0\), \(c\) is a fixed
source-invariant left multiplier. It can be moved from the input
to the end and removed from an image test by the same finite
Fourier--Gram argument as a centralizer multiplier. Polarization
shows that every such image effect commutes with \(c\). If a
bounded operator \(Q\) commutes with \(N_0\), then, for positive
\(c\in N_0\),
\[
 \langle Qf,cQf\rangle
   \leq\|Q\|^2\langle f,cf\rangle.
\]
The characterization of conditional expectation by these tests
implies
\[
 E_{N_0}\bigl((Qf)(Qf)^*\bigr)
        \leq\|Q\|^2E_{N_0}(ff^*).
\]
This applies to the generated algebra and its support projections,
and also after any fixed finite matrix amplification.

If a closed subspace is invariant under \(N_0\), the spectral
projections \(p_n=1_{[0,n]}(E_{N_0}(ff^*))\) preserve it. The
columns \(p_nf\) are relatively bounded and converge to \(f\),
because
\[
 \|(1-p_n)f\|_2^2
  =d\tau_{\mathcal N}
        \bigl((1-p_n)E_{N_0}(ff^*)\bigr)\longrightarrow0.
\]
For a relatively bounded \(f\), ordinary left spectral truncation
\(f_n=1_{[0,n]}(ff^*)f\) gives operator-bounded columns converging
in Hilbert norm, with
\(E_{N_0}(f_nf_n^*)\leq E_{N_0}(ff^*)\).

These approximants converge in the link form against a fixed
Hilbert \(b\). Indeed first replace \(b\) by a bounded \(b_l\).
For that bounded input, \eqref{eq:links-relative-form} bounds the
error by \(d\|b_l\|_\infty^2\|f_n-f\|_2^2\).
The tails \(b-b_l\) are uniformly controlled by
\eqref{eq:links-relative-bound}, since
\[
 E_{N_0}\bigl((f_n-f)(f_n-f)^*\bigr)
       \leq 4E_{N_0}(ff^*).
\]
Take \(n\to\infty\) first and then \(l\to\infty\).
This argument, and its version with the two inputs reversed,
extends all the bounded-input Gram identities to the stated
domains.

Finally, suppose a gap event annihilates every one-column input.
For a bounded end column in the two-column link, first approximate
its left multiplier by a finite Fourier sum. Each fixed shift
preserves the gap event in the limit, so each resulting term
vanishes. The approximation error is estimated on the uncut
bounded columns by Lemma~\ref{lem:links-gram-transport}. Taking
the approximation limit proves the same inclusion for bounded
pairs, with no use of a mixed tensor decomposition.
\end{proof}

\begin{lemma}[Additional invariant limiting variables]
\label{lem:links-invariant-variables}
Adjoin to the compact spectral data any bounded limiting tests which
are asymptotically unchanged by fixed left and right Fourier shifts.
The corresponding limiting effects retain base decomposability and
the measured equivariance established above. If a one-column effect
is also source face-parabolic invariant, including its inactive
simple factors, and its event or test is invariant under the
measured target shifts, that effect is scalar. For a mixed
two-ended effect, independent target-shift invariance and
independent contracted-radical invariance make its mass scalar
on both inputs. These assertions apply to invariant Borel events
of the limiting variables.
\end{lemma}

\begin{proof}
For a fixed Fourier shift, the assumed asymptotic identity of
bounded tests gives precisely the limiting intertwining identity
used in Lemma~\ref{lem:links-gram-transport}. Its finite-Fourier
approximation proof is unchanged. Base decomposability and
commutation with companions follow before taking limits. An
invariant one-column test therefore gives an element of
\(\mathcal N\). The source face-parabolic invariance, followed
by the active-factor radical argument and the inactive-factor
centralizer argument in Proposition~\ref{prop:links-link}, puts
this element in \(\mathcal N^G=\mathbb C\).

For mixed links, independent target invariance puts the respective
slices in \(\overline{\mathcal N}\) and \(\mathcal N\).
Independent contracted-radical invariance makes each slice scalar;
polarization and normal slice separation give the scalar joint
effect. This reasoning works for a submeasurement with additional
spectral slots when its indicated event has these invariances.

For example, fixed left and right multiplications change Cartan
coordinates by bounded amounts. If a radial variable is divided
by a scale tending to infinity, bounded continuous tests in its
chosen compactification are asymptotically unchanged whenever this
bounded displacement tends uniformly to zero in that compactified
coordinate. For the measurements obtained by adjoining such
rescaled radial coordinates to the original face data, the needed
source face-parabolic invariance also follows from the construction.
A contracted radical preserves the entire enlarged data law: after
sliding it to the far end, the near-identity transporter changes
the uncut bounded input by a Hilbert-small error, uniformly against
all bounded diagonal tests. A fixed centralizer element, including
an element of an inactive simple factor, can be slid to that same
end. Its base reindexing is removed by invariant probability, and
its left multiplier is removed by the finite-Fourier Gram argument.
The additional radial test causes no extra term in this argument,
because each fixed Fourier shift changes its limiting coordinate
by zero. Thus image-and-stratum tests together with these radial
coordinates are automatically face-parabolic invariant. In
particular their target-invariant masses are scalar even on a
pure vertex ray. This conclusion does not require that the
nonmixed link norm factor through scalars.

The argument supplies the required invariant limiting tests; no
derivative estimate for an arbitrary finite-scale discontinuous
mask is asserted. Establish the continuous determining identities
first. Their spectral measures then transport identically, giving
the assertion for invariant Borel events of the limiting
variables.
\end{proof}

\subsection{Ordering and genuine source actions}

There are two separate kinds of spectral operation. A single
limiting compact variable has its usual Borel functional calculus.
For several possibly noncommuting compressed effects, an ordered
integral initially means a product of single-slot continuous tests.
Smooth joint functions on compact real charts can be inserted by
absolutely summable separated expansions, and fixed-level
finite-place functions by finite separated partitions. The
coefficient bounds make these operations converge in norm. Support
of an ordered product on a closed relation means that it vanishes
on open rectangles disjoint from that relation; support alone does
not supply a general ideal of discontinuous joint functions.

On a buffered transverse projective box, normalized matrix
multiplication has denominators bounded below. Over a finite local
field choose ratio charts with a maximal coordinate. Matrix
operations on these compact charts are equicontinuous, so a clopen
test at a fixed level pulls back to tests at one uniformly bounded
level. Units and choices of a projective representative do not
alter this conclusion. At real places the corresponding maps have
bounded derivatives of each fixed order. The geometric estimates
for the variable shifts that occur in transverse products are
proved in Section~\ref{sec:transfer}; bounded shifts there are
counted in a compact subset of the \emph{whole} ambient product,
whose intersection with the lattice is finite.

The order of factors also matters in incidence arguments. Homogeneous
local sections let one shrink a violation neighborhood so that its
nearby flags admit the required incident values in fixed test
patches. Off a closed position relation, it is enough first to
prove vanishing on such rectangles. For a graph substitution in
totally disconnected slots, refine to finite clopen partitions
and insert the projection of the shared panel at the two
consecutive slots. Their locally constant factors then agree
exactly. At real slots, the corresponding substitution uses the
regular smooth graph equations established in the transfer
argument. This prescription retains the stated product order
throughout.

\begin{lemma}[Cancelling source multipliers]
\label{lem:links-genuine-source}
After establishing the decomposable link forms, one may enlarge
their source actions to genuine strongly continuous group
representations by an opposite-multiplier regular factor at each
used input base. The compressed effects become their pointwise
normal ampliations. This enlargement imposes no new mixing or
scalar-commutant assertion.
\end{lemma}

\begin{proof}
Write the projective source multiplier as \(\eta\). The regular
representation with multiplier \(\overline\eta\), induced with
the same section, has the conjugate of the scalar defect of
\(X_a(y)=\pi_0(c(a,y))\). The product of these two fibre
transporters consequently satisfies the exact group composition
law. The section continuity-in-measure argument gives strong
continuity of the resulting induced representation.

The added coordinate is unmeasured. At one base, its two
transported vectors have exactly their original inner product,
so a base-decomposable compressed form is tensored pointwise with
the identity. For a two-ended form the two base algebras commute
with its effect, and the same construction is performed
independently at the two inputs. The generator identification of
minimal spectral dilations gives the corresponding enlarged
dilation. In particular a face spectral measure already invariant
under the projective transport now commutes with a genuine source
representation and can be disintegrated with that representation.
All mixing and scalarity statements continue to refer to the
original finite algebra \(\mathcal N\); none is transferred to
the enlarged commutant.
\end{proof}

\subsection{Escape from bounded target coordinates}

\begin{lemma}[No all-bounded stratum]
\label{lem:links-no-bounded-stratum}
On every escaping source face path, the all-bounded gap stratum
has zero mass in a single measured lattice copy. This holds for
the one-column measurement and for bounded pairs in the link.
\end{lemma}

\begin{proof}
We first record the finite-module Bessel estimate used here. For
a bounded column \(b\in eL^2(B)^k\) and a Hilbert column \(\xi\),
Parseval in the regular coordinates gives
\[
 \sum_{h\in R}|\langle b w_h^*,\xi\rangle|^2
       =\|b^*\xi\|_2^2
       \leq\|b\|_\infty^2\|\xi\|_2^2.
\]
The harmless projective phases do not change the squared absolute
values. The same bound holds for a subset of translates, in
particular the \(R'_j\)-translates. In finite matrix versions one
sums this identity over the fixed row and column indices.

For bounded end and start columns, the squared norm of a coefficient
at a fixed measured index \(g\in\Gamma_j\) in
\(b(ay)^*X_a(y)f(y)\) is the sum over \(R'_j\) of squared
frame coefficients of the pair. Equivalently these are inner
products of \(X_a(y)f(y)\) with the corresponding partial
translates of \(b(ay)\), with \(g\) absorbed into one of the
columns. This is just Parseval in the remaining group coordinates,
and remains true with their twisted coefficient phases.

Use now the complementary regular realization for
\(\Gamma_0\times R'_j\). Project each column to a finite
\(\Gamma_0\)-coordinate window. These projections commute with
the partial \(R'_j\)-translates, so the Bessel bounds are
preserved. Source transport is unitary and commutes with those
partial translates, and the fixed measured shift likewise
preserves their bounds. Both projected columns converge in
integrated Hilbert norm to the original columns as the finite
windows increase. The Bessel estimate on the other column bounds
the sum of squared coefficient errors by that Hilbert error.
Interchanging the two columns gives the estimate for the second
error. Thus the finite-window approximation is uniform in the
path parameter. Reindexing the end column at \(ay\) does not
change its averaged Hilbert error.

For fixed finite windows, the translated supports are disjoint
unless the source cocycle \(c(a,y)\) belongs to a fixed finite
subset of \(\Gamma_0\). That exceptional event has measure
tending to zero. Indeed, for any \(\delta>0\), choose a compact
set containing the section values \(d_0(y)\) outside a set of
measure \(\delta\). Invariance of quotient probability gives
the same assertion for \(d_0(ay)\). On the intersection of the
two good sets, if \(c(a,y)\) belongs to a fixed finite set
\(L\), then
\[
 a=d_0(ay)c(a,y)d_0(y)^{-1}
\]
belongs to one fixed compact subset of \(G\). This is impossible
eventually on an escaping path. The exceptional measure is at
most \(2\delta\), and \(\delta\) is arbitrary.

Off that exceptional set, every partial-translate coefficient of
the finite-window pair is zero. On it, their summed squares are
uniformly bounded by the Bessel estimate and the column bounds.
Integration gives convergence to zero. Remove the finite-window
approximation afterward. We have proved vanishing of the norm
at every fixed measured \(\Gamma_j\)-index.

A bounded Cartan set is compact in the whole ambient product, so
contains only finitely many elements of the discrete lattice.
Summing the preceding conclusion gives zero link mass on each
bounded Cartan window. Increasing those windows removes the
all-bounded stratum. This is a count in the whole lattice and
requires no discreteness of a projection to a proper subproduct.
Finally apply \eqref{eq:links-frame}, with its finite corner
frame, to obtain the one-column assertion.
\end{proof}

\subsection{An amenable-boundary obstruction}

The measured target boundary in the following argument is the full flag
space of the copy \(\Gamma_j\). The arbitrary comparison group
\(\Lambda\) is among the commuting companion actions.

\begin{lemma}[Full-flag obstruction]
\label{lem:links-obstruction}
Let \(S\) be the root-generated group of a noncompact isotropic
almost simple component of a semisimple source Levi subgroup, using
its root-generated lift from Convention~\ref{conv:setup-panels} when
the ambient action requires it. There is no nonzero positive map from the
continuous functions on the full target flag space to
\(\mathcal A_j\) which is target-equivariant and fixed by the source
action \(\sigma|_S\).
\end{lemma}

\begin{proof}
Write \(Z=G/P\) for the full target flag space. Suppose that
\(\Phi:C(Z)\to\mathcal A_j\) is such a map. Its mass
\(a=\Phi(1)\) belongs to \(\mathcal N^S\): target equivariance makes
it commute with the measured target action, and membership in
\(\mathcal A_j\) gives the remaining commutations. Choose
\(\epsilon>0\) such that
\(p=1_{[\epsilon,\infty)}(a)\ne0\). On the reducing space
\(p\mathcal H\), the formula
\[
 \Psi(u)=p a^{-1/2}\Phi(u)a^{-1/2}p
\]
defines a unital completely positive map with the same covariances.
A positive map from the commutative algebra \(C(Z)\) is completely
positive. The displayed normalization does not require \(a\) to
commute with \(\Phi(u)\).

Choose a discrete free subgroup \(F\cong\mathbb F_2\) in \(S\).
One obtains it from two sufficiently strong rank-one contractions with
disjoint attracting and repelling neighborhoods; at a finite place the
same construction uses translations in the rank-one tree. Passing
through a finite central cover leaves a free subgroup after choosing
the generators in that cover. Let \(\mathcal D\) be the concrete
\(C^*\)-algebra on \(p\mathcal H\) generated by the cut base
multipliers, the companion actions \(W|_{R'_j}\), and
\(U_f\), \(f\in F\). These operators commute with both the measured
target action and every value of \(\Psi\).

We first obtain a completely contractive minimum-tensor-norm estimate.
We use the minimal Stinespring construction and its commutant lifting
\cite{Stinespring1955} and \cite[Theorems~1.1.1 and~1.3.1]{Arveson1969}.
In the minimal construction of \(\Psi\), the algebraic
vectors are \(u\otimes\xi\), with form
\[
 \left\langle\sum_i u_i\otimes\xi_i,
                  \sum_l v_l\otimes\eta_l\right\rangle
 =\sum_{i,l}\langle\xi_i,\Psi(\overline{u_i}v_l)\eta_l\rangle.
\]
The lift of \(D\in\mathcal D\) is
\(\widehat D(u\otimes\xi)=u\otimes D\xi\).
Commutation with the matrix of values of \(\Psi\) gives
\(\|\widehat D\|\leq\|D\|\), with the concrete input norm on
the right. Thus the companion representation is controlled before
dilation; no larger norm on a newly presented companion algebra is
introduced. If \(W_gW_h=\zeta(g,h)W_{gh}\) in the measured copy,
the covariant lift
\[
 T_g(u\otimes\xi)=(g\cdot u)\otimes W_g\xi
\]
has the same multiplier \(\zeta\).

The action of \(\Gamma_j\) on \(Z\) is topologically amenable.
By \cite[Proposition~2.5]{AnantharamanDelaroche2002}, choose continuous
fields \(\eta_i:Z\to\ell^2(\Gamma_j)\) with finite coordinate support,
whose norms tend to one and whose equivariance errors tend to zero
uniformly on \(Z\) for each fixed \(g\in\Gamma_j\).
Compactness of \(Z\) makes these norms uniformly positive eventually;
normalizing pointwise gives continuous unit fields
\(\xi_i(z)=\eta_i(z)/\|\eta_i(z)\|\) with the same uniform
asymptotic equivariance. If \(\pi_Z\) is the
spectral representation of \(C(Z)\) in the dilation, set
\[
 Q_i\eta=\sum_{t\in\Gamma_j}\delta_t\otimes
                    \pi_Z(\xi_i(\cdot)_t)\eta.
\]
The maps \(Q_i\) are isometries. They intertwine the companion
algebra exactly, and
\(Q_iT_g-(\lambda_g\otimes T_g)Q_i\) tends to zero in operator
norm for each fixed \(g\). Projective absorption is explicit:
the unitary
\[
 A(\delta_t\otimes\eta)=\delta_t\otimes T_t^*\eta
\]
conjugates \(\lambda_g\otimes T_g\) to
\(\lambda_{\zeta}(g)\otimes1\). It leaves the companion
representation on the second factor unchanged. Taking limits and
restricting to the original input proves, at every finite matrix level,
\begin{equation}
 \left\|\sum_g W_gD_g\right\|
 \leq
 \left\|\sum_g\lambda_{\zeta}(g)\otimes D_g\right\|_{
                  C^*_{r,\zeta}(\Gamma_j)\otimes_{\min}\mathcal D}.
 \label{eq:links-min-bound}
\end{equation}
This also fixes the multiplier convention: the multiplier in this
formula is that of \(W\), and is conjugated if one instead writes the
right action with the adjoint generators of \(B\).

Complementary regularity makes
\(H\) a subrepresentation of a finite regular multiple for
\(\Gamma_0\times R'_j\), with split multiplier. On inducing the
\(\Gamma_0\) coordinate, the identification
\[
 (y,h)\longmapsto d_0(y)h
\]
turns it into the \(G\) coordinate. In these coordinates \(U_f\)
acts as left translation on \(G\), times a scalar function of modulus
one; the projective regular multiplication only contributes that
scalar. The companions \(W|_{R'_j}\) use a separate twisted regular
coordinate. Tensoring with the separate measured regular
representation in \eqref{eq:links-min-bound}, a simple-base finite
Fourier matrix field from the full split group algebra of \(R\) acts
fibrewise with its original matrix norm.

Consequently, for contraction fields \(C_f\), scalar phases
\(\theta_f\), and a finitely supported probability \(\nu\) on
\(F\), the operator in this spatial model satisfies
\[
 \left\|\sum_{f\in F}\nu(f)C_f\theta_f L_f\right\|
       \leq\|\lambda_F(\nu)\|.
\]
Indeed, pointwise Hilbert norms are bounded by
\(\sum_f\nu(f)\|\eta(f^{-1}g)\|\). Disintegrating Haar measure
along the left \(F\)-orbits in \(G\) gives the ordinary positive
regular convolution bound on each orbit. There is no commutation
assumption between \(C_f\) and \(L_f\). For the concrete cut
representation this reasoning is applied to the indicated generator
polynomials and then restricted to their invariant input subspace; it
does not require an extension of every abstract element of the cut
companion algebra to the enlarged regular space.

We contradict this estimate on one matrix vector. Regard the matrix
\(E(y)=p(y)\) as \(k\) columns, and put, with \(x=f^{-1}y\),
\[
 T_f(y)=p(x)X_f(x)^*p(y).
\]
This is a contraction. Since \(p\) is \(\sigma_f\)-invariant,
\(X_f(x)p(x)X_f(x)^*=p(y)\), and right multiplication after the
shift gives
\[
 (R_{T_f}U_f^{(k)}E)(y)
   =X_f(x)p(x)\,p(x)X_f(x)^*p(y)=p(y).
\]
Right matrix multiplication has the opposite convention
\((R_T)_{\beta\alpha}=R_{T_{\alpha\beta}}\); the transpose of
matrix indices implements
\(M_k(B)^{\mathrm{op}}\cong M_k(B^{\mathrm{op}})\).
In particular it preserves the contraction norm required by the
spatial estimate.

For the finitely many \(f\) in the support of \(\nu\), choose
contraction fields \(T_{f,n}\), simple in the base and finite in
Fourier support, with \(T_{f,n}\to T_f\) in integrated tracial
\(L^2\). To see the contractive approximation directly, first
approximate in the unit ball by simple-base fields, use Kaplansky
density of the reduced group algebra in its group von Neumann algebra,
then approximate each reduced-algebra entry in operator norm by a
Fourier polynomial and rescale by its at most \(1+o(1)\) norm.
On the bounded matrix vector just used,
\[
 \|(R_{T_{f,n}}-R_{T_f})U_f^{(k)}E\|_2
       \leq\|T_{f,n}-T_f\|_2.
\]
Thus \(S_n=\sum_f\nu(f)R_{T_{f,n}}U_f^{(k)}\) satisfies
\(S_nE\to E\ne0\). On the other hand,
\eqref{eq:links-min-bound} and the spatial estimate give
\(\|S_n\|\leq\|\lambda_F(\nu)\|\).
For the uniform walk on two free generators and their inverses this
norm is at most \(\sqrt3/2<1\), the classical free-group estimate of \cite[Theorem~3]{Kesten1959}. For the needed upper bound, the weight
\(3^{-|h|/2}\) on the four-regular tree gives this bound by the
weighted Schur estimate. The two conclusions are incompatible because
\(0<\|E\|_2^2=\mathsf T(p)<\infty\).
\end{proof}

\subsection{Strict growth under a slower Levi translation}

The next argument follows one positive stratum through an outer
limit. It uses invariance of its image marginal, but does not
assume invariance of its reverse data or residual projectivity.

\begin{lemma}[A slower Levi translation removes every compact window]
\label{lem:links-strict-growth}
Let an inner source face path have a positive target gap stratum
\(\Omega_D\). After the inner limit, append a slower commuting
split translation whose projection escapes in a noncompact simple
component \(S\) of the remaining source Levi. Additional split
parts commuting with \(S\) are allowed. Every fixed compact
window in \(\Omega_D\) has mass tending to zero in the outer
limit. The entire positive contribution of \(\Omega_D\) then
passes to strata \(\Omega_{D'}\) with \(D'\supsetneq D\).
\end{lemma}

\begin{proof}
Its scalar mass is some \(m_D>0\), by
Proposition~\ref{prop:links-link}. Divide the stratum measurement
by \(m_D\), and denote the resulting unital full-data measurement
by \(\Psi\). If \(D=\Pi\), its image measurement is a
nonzero full-flag map fixed by the remaining source component
\(S\), contradicting Lemma~\ref{lem:links-obstruction}.
Thus in the situation under consideration \(D\ne\Pi\).

\emph{The marginal inclusion.}
Let
\(\Psi(F)=V_\Omega^*M_FV_\Omega\) be a minimal full-data
dilation, and let \(x:\Omega_D\to X_D\) be its image-flag map.
Write \(V:\mathcal H\to\mathcal K_X\) for the minimal image
dilation. The generator assignment
\[
 M_uVf\longmapsto M_{u\circ x}V_\Omega f
\]
preserves inner products. It extends to an isometric inclusion
\(i:\mathcal K_X\to\mathcal K_\Omega\), with
\(iV=V_\Omega\). This inclusion intertwines the image spectral
measure, base multiplication, all companions, and the canonical
target lifts. Disintegrating over the image variable gives
\(\mathcal K_X=\int^\oplus K_x\,d\vartheta(x)\) and
isometric fibre inclusions \(i_x\). The base variables and their
commuting decompositions are retained in this notation.

\emph{Positive density forms.}
Let \(\mathcal F_x\) be the compact full-refinement space of the
image face \(x\), with compact-rotation probability \(\eta_x\).
The residual projectivity of a full datum \(\omega\in\Omega_D\)
transports the compact-rotation probability on its reverse
refinement space to a probability on \(\mathcal F_x\). Write
this probability as
\[
 k_\omega(l)\,d\eta_x(l),\qquad
 k_\omega(l)>0,\qquad
 \int_{\mathcal F_x}k_\omega(l)\,d\eta_x(l)=1.
\]
On a compact data window \(C\subset\Omega_D\), the density
has bounds \(0<c_C\leq k_\omega(l)\leq C_C<\infty\), uniformly
in \(\omega\in C\) and \(l\in\mathcal F_x\).
Indeed the residual transformations and their inverses range in
a compact set on such a window. Their Jacobians against the
compact-rotation measures are positive and uniformly bounded.
At a finite place the same conclusion uses the analytic Haar
densities in finitely many compact-open charts and bounded
absolute Jacobians. No real differential operator is used there.

For fixed \((x,l)\), define bounded increasing operators
\[
 D_{n,x}(l)=i_x^*M_{\min(k_\omega(l),n)}i_x.
\]
Their limiting quadratic form is
\[
 q_{x,l}(v)=\|M_{\sqrt{k_\omega(l)}}i_xv\|^2.
\]
This form is closed: multiplication by
\(\sqrt{k_\omega(l)}\) is closed and its composition with the
isometric inclusion \(i_x\) is closed. Its domain is dense for
almost every \((x,l)\). To see this, choose countably many
input vectors \(f_r\) such that the sections \(Vf_r(x)\) are
total in \(K_x\) almost everywhere. Indeed, the measurable projection
onto their fibrewise orthogonal complement kills every angular
multiple of \(V\mathcal H\); minimality of the image dilation makes
that projection zero.
Tonelli and the normalization of \(k\) give
\[
 \int_{X_D}\int_{\mathcal F_x}
        q_{x,l}(Vf_r(x))\,d\eta_x(l)d\vartheta(x)
          =\|f_r\|^2.
\]
Thus their span lies in the form domain almost everywhere, on
one common conull set. Let \(D_x(l)\) be the associated positive
self-adjoint operator. The bounded measurable approximating
forms above, equivalently their resolvents, give its measurable
dependence. For every \(v\in K_x\), integration of the
nonnegative form yields
\[
 \int_{\mathcal F_x}q_{x,l}(v)\,d\eta_x(l)=\|v\|^2.
\]
Consequently
\begin{equation}
 (Tf)(x,l)=D_x(l)^{1/2}(Vf)(x)
 \label{eq:links-density-isometry}
\end{equation}
is a bounded isometry from \(\mathcal H\) to
\[
 \mathcal R=
 \int_{X_D}^{\oplus}L^2(\mathcal F_x,\eta_x;K_x)\,d\vartheta(x).
\]

\emph{Equivariance and source lifts.}
For a target element \(\gamma\), set
\[
 r_{\gamma,x}(l')=
       \frac{d(\gamma_*\eta_x)}{d\eta_{\gamma x}}(l').
\]
Transport of the residual probability gives
\(k_{\gamma\omega}(\gamma l)
  =r_{\gamma,x}(\gamma l)k_\omega(l)\).
Since \(i\) intertwines the canonical target lifts, the
compressed density forms have the same covariance. If
\(u_{\gamma,x}:K_x\to K_{\gamma x}\) is the fibre-unitary
part of the image lift, their operator form is
\[
 D_{\gamma x}(\gamma l)
   =r_{\gamma,x}(\gamma l)\,
            u_{\gamma,x}D_x(l)u_{\gamma,x}^*.
\]
The target lift on \(\mathcal R\) includes the half-density
\(r_{\gamma,x}^{1/2}\), in addition to any factor already
present in the image direct integral. Taking square roots proves
that \(T\) intertwines the target action. Base and companion
operators commute with full-data multiplication and intertwine
\(i\), so \(T\) intertwines them as well. Thus \(T\)
intertwines base multiplication and all of \(W_R\).

For a source centralizer element \(s\), invariance of the image
marginal gives a canonical lift on its minimal dilation, defined
on generators by
\[
 U_s^X(M_uVf)=M_uVU_sf.
\]
It fixes the image coordinate and is strongly continuous, first
on these generators and then by density. Extend it to
\(\widetilde U_s\) on \(\mathcal R\) by its action on the
\(K_x\) fibre, with the induced base reindexing, and the identity
on \(l\). It commutes with multiplication by the full
refinement flag. These lifts use only the image marginal; no
lift preserving the full-data law is being assumed.

\emph{The finite-trace intertwiner space.}
Let \(\mathscr E_0\) be the bounded intertwiners from
\(\mathcal H\) to \(\mathcal R\) for base multiplication and
all target actions \(W_R\). For \(B,C\in\mathscr E_0\),
\(B^*C\in\mathcal N\), so define
\[
 \langle B,C\rangle_{\mathscr E}
       =\tau_{\mathcal N}(B^*C)
\]
and take the Hilbert completion \(\mathscr E\). There are two
useful exact finite-trace estimates:
\begin{equation}
 \|B\|_{\mathscr E}^2
       =\frac1d\sum_{r=1}^k\|Bb_r\|^2,
 \qquad
 \|Bf\|^2\leq d\|f\|_\infty^2\|B\|_{\mathscr E}^2
 \quad(f\text{ operator-bounded}).
 \label{eq:links-intertwiner-frame}
\end{equation}
Both follow by testing \(B^*B\) against the corner frame, or
against the bounded positive matrix \(ff^*\), respectively.

Source conjugation
\(\mathscr U_sB=\widetilde U_sBU_s^*\) is isometric for this
norm, because \(\sigma_s\) preserves
\(\tau_{\mathcal N}\). Its base-valued projective defects
cancel since \(B\) intertwines the base algebra. Hence it is a
genuine unitary representation on \(\mathscr E\).
Strong continuity follows from the frame equality in
\eqref{eq:links-intertwiner-frame}, strong continuity of the two
source lifts, and boundedness of \(B\), then extends by density
to the completion.

Here \(S\) is the root-generated Levi group, or its root-generated
lift from Convention~\ref{conv:setup-panels}; invariant spaces are
taken for that group.  This representation has no \(S\)-invariant
vector. If it did,
choose a bounded intertwiner \(B\) whose orthogonal projection
onto the invariant subspace is nonzero. The closed convex hull
of its orbit contains that projection. Finite convex combinations
\(B_n\) therefore converge to it in \(\mathscr E\), and
\(\|B_n\|_{\mathrm{op}}\leq\|B\|_{\mathrm{op}}\).
The bounded-column estimate in
\eqref{eq:links-intertwiner-frame} gives strong convergence on
bounded columns. Their density and the common operator bound
give a bounded limiting operator \(B_\infty\). It still
intertwines the base and target actions, represents the nonzero
invariant vector, and is \(S\)-fixed. The map
\[
 u\longmapsto B_\infty^*M_{u(l)}B_\infty
\]
is then a nonzero positive full-flag map into \(\mathcal A_j\),
target-equivariant and \(S\)-invariant. This contradicts
Lemma~\ref{lem:links-obstruction}.

Howe--Moore now gives weak decay of translates of \(T\) along
escape in \(S\), also when multiplied by the additional
commuting translations. For precision, disintegrate the
representation into type-I irreducibles of \(S\) with
multiplicity. There is no trivial summand. Commuting unitaries
act on the multiplicity spaces. Approximation by finite tensor
sums reduces the assertion to finitely many matrix coefficients
of the irreducible \(S\)-representations, which tend to zero
uniformly over the multiplicity unitaries. Dominated convergence
finishes the argument.

\emph{The positive compact-window comparison.}
For a compact window \(C\subset\Omega_D\), let
\(L_C=i^*M_{1_C}i\). This is a decomposable positive contraction
on \(\mathcal K_X\), and
\(V^*L_CV=\Psi(C)\). Its density lower bound gives
\(D_x(l)\geq c_C L_C(x)\) in form order. Operator monotonicity
of the square root, followed by
\(L_C^{1/2}\geq L_C\), yields
\begin{equation}
 D_x(l)^{1/2}\geq\sqrt{c_C}\,L_C(x).
 \label{eq:links-positive-window}
\end{equation}
One may obtain this inequality directly from the bounded
\(D_{n,x}(l)\) for \(n\geq c_C\) and then pass to their
closed-form limit.

Write the extended path as \(a_t b_s\), with the inner
\(t\)-limit first. The exact cocycle formula and the change of
base variable \(z=b_sy\) show that the inner full-data law is
precomposed by \(U_{b_s}\):
\[
 \Psi_s(F)=U_{b_s}^*\Psi(F)U_{b_s}.
\]
Its image marginal is unchanged. The appropriate translated
intertwiner is therefore
\(T^{(s)}=\widetilde U_{b_s}^*TU_{b_s}\).
Because the marginal lift fixes \(x,l\), its fibre expression
is the positive operator
\((U_{b_s}^X)^*D_x(l)^{1/2}U_{b_s}^X\) applied to \(Vf\),
with base reindexing understood. Inequality
\eqref{eq:links-positive-window} remains valid after this
conjugation and bounds the correspondingly translated \(L_C\).

Put \((J_0f)(x,l)=Vf(x)\), the constant-refinement section.
For bounded \(f\), the functional
\(B\mapsto\langle Bf,J_0f\rangle\) is continuous on
\(\mathscr E\), by \eqref{eq:links-intertwiner-frame}.
It is unnecessary that \(J_0\) itself intertwine the target
action. Weak decay and the positive inequality give
\[
 0\leq\sqrt{c_C}\,\langle f,\Psi_s(C)f\rangle
       \leq\langle T^{(s)}f,J_0f\rangle\longrightarrow0.
\]
The right side is real and nonnegative by its positive fibre
expression. Density of bounded columns and
\(0\leq\Psi_s(C)\leq1\) extend the vanishing to every input.
This proves disappearance of the actual translated window mass,
rather than merely weak convergence of an unrelated auxiliary
operator.

\emph{The entire stratum contribution.}
Each precomposed law \(\Psi_s\) has mass one and takes values on
\(\Omega_D\). Take its outer limit on
\(\overline{\Omega_D}\). Every continuous test compactly
supported in \(\Omega_D\) has zero limiting effect by the
window estimate. A countable compact-window exhaustion therefore
shows that the outer law gives no mass to \(\Omega_D\).
Lemma~\ref{lem:links-matrix-data} places its remaining support in
\[
 \overline{\Omega_D}\setminus\Omega_D
       \subseteq\bigcup_{D'\supsetneq D}\Omega_{D'}.
\]
Restoring the scalar factor \(m_D\) proves the assertion for
the entire original positive contribution. There are only
finitely many gap sets, so the same outer limit tracks all of
their contributions. Further subpieces are dominated by these
positive total measurements. At the next stage the covariance
and scalar-mass statements apply anew to each positive total
stratum of the enlarged source face; invariance of the original
full-data law was never needed.
\end{proof}

\begin{proposition}[Gap count and vertex quotients]
\label{prop:links-gap-count}
In one measured copy, every positive stratum of a source path
toward a face of type \(I\) has at most \(|I|\) infinite gaps.
A vertex path has exactly one infinite gap. In a successive
single-gap lex chain starting at a vertex, the stage with
\(k\) source gaps has exactly \(k\) target gaps; in particular
each such extension adds exactly one.

For a quotient of a faster and a slower vertex translation in a
common apartment, there is exactly one target gap if the two
source vertices are coaxial, and at least two if they are not.
Here coaxial means equal or opposite vertices on the same
translation axis.
\end{proposition}

\begin{proof}
Put \(r=|\Pi|\), and start with a positive target stratum \(D\)
on a source face of type \(I\). Append the \(r-|I|\) missing
source coordinates at successively slower scales. At each stage
the fundamental split vector for a missing root has a nonzero
projection to the remaining-Levi component containing that root;
its other split parts commute with that component. Thus
Lemma~\ref{lem:links-strict-growth} applies at each stage. If a
full target stratum were reached before the source face were
complete, the obstruction in that lemma would already give a
contradiction. Otherwise the tracked positive contribution must
pass through a strictly increasing chain of gap sets of length
\(r-|I|\). Since a target has only \(r\) simple types,
\[
 |D|+r-|I|\leq r,
\]
which proves \(|D|\leq|I|\). If \(I=\Pi\), this is just the
automatic ceiling \(|D|\leq r\), and no remaining-Levi
obstruction is invoked.

For a vertex, Lemma~\ref{lem:links-no-bounded-stratum} excludes
zero gaps and the upper bound excludes two or more. The same
strict-growth argument, applied successively from that vertex,
gives at least \(k\) gaps at stage \(k\), whereas the upper
bound gives at most \(k\). This proves the exact lex counts.
More generally, an extension from a stage already attaining its
upper bound has an exact increment of one. No exact-increment
claim from an arbitrary unsaturated simultaneous multi-gap
initial path is needed.

Finally choose the common apartment for the two vertex
translations and make the faster vertex a dominant fundamental
ray \(v\), of type \(j\). If the slower vertex is coaxial,
their quotient is again a vertex translation at the faster
scale, possibly with the opposed orientation, so the vertex
conclusion applies. If its vector is \(w\) and is not coaxial,
then it has a nonzero projection to a simple factor of the
remaining Levi. Indeed the common kernel of the remaining simple
roots \(\Pi\setminus\{j\}\) is precisely \(\mathbb Rv\).
Thus a zero projection would force \(w\) to lie on that line,
which is the coaxial case. After choosing a dominant chamber
inside the remaining Levi, the slower quotient component
therefore escapes in a noncompact simple factor; the other
split components commute with it. Strict growth from the
one-gap faster limit forces at least one additional target gap.
All these statements retain the faster-first order of limits.
\end{proof}
